\subsection{计算分数选择的严格 CPU 参考}
\label{sec:7-strict-selection}

补证 T4 使用 GPU 生成的 FP32 分数，在 CPU 上按
$(\primarydist,\mathrm{position})$ 对每条查询的全部已访问候选排序。
诊断导出器遵循打包扫描布局和累加顺序；所有返回有效主码项的导出分数与选择器
分数逐位一致。源码对照确认，诊断版本相对拆表实验实现只改变导出部分。
Vogue、arXiv、OpenAI-3072 采用 1/8/128 探针、
$\alpha\in\{4,100,200\}$、$\gscan\in\{0,1,2,4,8\}$ 及共享界开／关，
共 270 个配置，每个配置 256 条查询，即 69,120 个查询实例
（\cref{tab:7-strict-selection}）。

\suppressfloats[t]
\begin{table}[htbp]
\centering
\caption[严格计算分数选择检查]{T4 CPU 参考检查。集合及顺序一致使用完整主码键。
有效数量检查 $\min(c,n_c)$ 个 survivor；末端一致性按
$(\refineddist,\mathrm{ID})$ 排序已保留的计算精排分数，包含 FP32 比特比较。
计数为查询实例，同一查询在配置间重复使用。}
\label{tab:7-strict-selection}
\footnotesize\setlength{\tabcolsep}{3pt}
\begin{tabular}{@{}lrrrrr@{}}
\toprule
路径 & 查询数 & 集合／顺序精确 & 纯并列差异 & 有效数量正确 & 末端精确 \\
\midrule
$\gscan\geq1$ & 55,296 & 55,296 / 55,296 & 0 & 55,296 & 55,296 \\
$\gscan=0$ & 13,824 & 13,768 / 8,758 & 56 & 13,824 & 13,824 \\
\bottomrule
\end{tabular}

\end{table}

55,296 个分组实例全部与完整键参考的顺序一致。
单 block 的 13,824 个实例全部满足距离全预序，56 个集合差异全部位于截断并列。
19,120 个候选不足实例均返回正确数量的有效 survivor；全部 69,120 个末端选择
均与保留分数参考一致。这直接检查
\cref{lem:5-buffer,prop:5-local-coverage,prop:5-shared-bound} 的阶段契约。
此前 64 个扫描对照及 48 个共享界对照仍提供独立的实现间配对证据；归档 sanitizer
覆盖较小子集。导出与归档边界见\cref{sec:a-t4-t5}。
